% 本文件由 LaTeX 排版 Agent 根据有效 translation.json 自动生成。
% author=Augustine of Hippo; work_id=1405; title=驳斥摩尼的基本书信

\chapter{驳斥\Person{摩尼}的基本书信}

% structure: p0001 source=h1 target=omitted reason=page-title-already-rendered-by-parent

% structure: p0002 source=h2 target=section reason=relative-heading-rank; base=h2

\section{医治异端者胜于消灭他们}

1. 我向惟一真实、全能的天主祈祷，万物都源于祂、藉着祂、并在祂内；我过去和现在都祈求祂，在反对和驳斥你们摩尼教徒的异端时（你们或许更多的是由于轻率而非恶意而成为异端者），赐给我一颗平静安详的心，旨在使你们归正而非击败你们。因为当主藉着祂的仆人推倒错误的国度时，祂对待迷途之人（就他们是人而言）的旨意是使他们被改正而非被消灭。并且在最终审判之前，每当天主施予惩罚时，无论是藉着恶人或义人，无论是藉着无知者或有知者，无论是隐秘地或公开地，我们都必须相信其设计的效果是医治人，而非毁灭他们；同时，对于那些拒绝康复方法的人，是在为最终结局作准备。因此，正如宇宙中有些事物用于肉身惩罚，如火、毒药、疾病等，另一些事物则用于惩罚心灵，不是通过肉身的痛苦，而是通过自身情欲的缠绕，如损失、流放、丧亲、辱骂等；而另一些事物，不带来折磨，却适合于安慰和安抚衰弱者，例如安慰、劝勉、讨论等；在这一切中，天主的至高正义有时利用甚至恶人（他们无知地行事），有时利用善人（他们有智慧地行事）。因此，我们应当更愿意选择更好的部分，即我们得以达成使你们改正的目的，不是通过争辩、争斗和迫害，而是通过仁慈的安慰、友善的劝勉、平静的讨论；正如经上所载：\BibleQuote{主的仆人不可争辩，但要温和对待众人，善于教导，忍耐；用温柔教导那些反对自己的人。}\BibleRef{弟后 2:24} 所以，我说，我们应当渴望在这项工作中获得这一部分；将善赐给那些渴望并祈求祂的人，是属于天主的。\par

% structure: p0004 source=h2 target=section reason=relative-heading-rank; base=h2

\section{为何应当更温和地对待摩尼教徒}

2. 让那些不知道寻找真理需要多大辛劳、避免错误需要多大困难的人向你们发怒吧。让那些不知道以虔诚心性的宁静克服肉身幻想是多么罕见和艰难的人向你们发怒吧。让那些不知道医治内心之人的眼睛使他可以注视他的太阳——不是你们所崇拜的那太阳，那个在血肉之人和野兽眼中以天体的光辉照耀的太阳——而是那一位，正如先知所记载的：\BibleQuote{正义的太阳在我身上升起，}\BibleRef{拉 4:2} 以及福音中所说的：\BibleQuote{那才是真光，照亮每一个来到世上的人。}\BibleRef{若 1:9} 让那些从未像他们看到你们那样被误导的人向你们发怒吧。\par

% structure: p0006 source=h2 target=section reason=relative-heading-rank; base=h2

\section{\Person{奥古斯丁}曾为摩尼教徒}

3. 至于我——在长期且持续的困惑之后，终于发现了单纯的真理，这真理无需记载于任何虚构的传奇即可习得；我，不幸的人，藉着天主的帮助，勉强驳斥了从各种理论和错误中收集而来的内心虚妄想象；我如此晚才寻求医治我心灵昏昧的方法，遵从了全慈悲的医生的召唤与温柔劝勉；我曾长久哭泣，恳求那恒常不变、不可侵犯的存在，按照圣经的见证，屈尊向我内心印证祂自己；藉着祂，最终，所有那些因你们长期熟悉而对你们有如此牢固掌控的一切虚构，我都曾勤奋审查、仔细听闻、过于轻易相信、在一切机会下向他人推荐相信、并坚决勇敢地反对对手——我绝不能向你们发怒；因为我必须容忍你们，如同从前我不得不容忍自己一样，并且我必须对你们有耐心，如同我的同伴们对我那样，当时我疯狂而盲目地迷失在你们的信仰中。\par

4. 另一方面，所有人都必须承认，作为回报，你们也应当放下所有的傲慢，这样你们才会更倾向于温和，而不会以敌对的精神反对我，导致自身的伤害。我们双方都不可断言自己已经找到了真理；让我们寻求它，仿佛它对我们都还是未知的。因为如果真理没有被任何轻率的预设认为已经找到并确定，就可以以热情和一致来寻求。但如果我无法说服你们同意这一点，至少允许我假设自己是一个陌生人，现在第一次听到你们，第一次审查你们的教义。我认为我的要求是公正的。并且必须作为一个公认的事实来规定：除非你们就所有涉及灵魂救恩的问题给予我清晰无模糊的解释，否则我不会与你们一起祈祷，或一起举行集会，或使用\Person{摩尼}的名号。\par

% structure: p0009 source=h2 target=section reason=relative-heading-rank; base=h2

\section{公教信仰的证明}

5. 因为在天主教会内——姑且不提最纯净的智慧，即少数属神的人在此生所能达到的认识，这种认识虽然极其有限（因为他们是人），却毫无不确定性（其余的大众之所以获得完全的安全感，不是靠才智的敏锐，而是靠信德的朴实）——姑且不提这种智慧（你并不相信它存在于天主教会内），还有许多其他事物最公义地将我留在她的怀抱中。万民万邦的共识使我留在教会中；同样，她的权威——由奇迹所开创，由望德所滋养，由爱德所扩大，由岁月所确立——也使我留在其中。司铎的传承使我留下，这传承从宗徒之长的伯多禄的宝座起——主在复活后曾委托他牧养自己的羊——一直延续到现今的主教团。最后，还有\Quote{天主教}这一名称本身，教会并非无理由地在众多异端中如此保留它；以至于，虽然所有异端者都希望被称为天主教徒，但当一位陌生人询问天主教会在哪里聚会时，没有哪个异端者敢于指出他自己的小教堂或房屋。因此，数目众多且重要的、属于基督徒之名的宝贵纽带，使信徒留在天主教会内，这是理所当然的，尽管由于我们理解力迟钝或我们生命成就微小，真理可能尚未完全显明。但在你们那里，没有任何这些事物来吸引或留住我，唯有真理的应许发挥作用。如果真理被证明得如此清晰，以至于没有任何怀疑的余地，那么它必须优先于所有使我留在天主教会内的事物；但如果只有应许而没有实现，没有人能将我从那以如此众多且强韧的纽带将我绑在基督宗教上的信仰中移走。\par

% structure: p0011 source=h2 target=section reason=relative-heading-rank; base=h2

\section{第五章——驳斥摩尼《书信》的标题}

6. 那么，让我们看看摩尼教导我什么；尤其要检查他称为根本书信的那篇论著，其中几乎包含了你们所信的一切。因为在那不幸的时期，我们阅读它时，按照你们的意见是得了光照。书信的开头是：——\Quote{摩尼，耶稣基督的宗徒，因天主圣父的眷顾。这些是从常存活泉涌出的有益言语。}现在，请你们耐心地听我追问。我不相信摩尼是基督的宗徒。我恳求你们，不要发怒并开始诅咒。因为你们知道，我的原则是不加考虑就不相信你们任何陈述。因此我问：这个摩尼是谁？你们会回答：基督的宗徒。我不相信。现在你们不知该说什么或做什么；因为你们应许传授真理的知识，而这里你们却强迫我相信我所不知道的事。也许你们会向我宣读福音，并试图在那里找到为摩尼作证的证据。但如果你们遇到一个还不相信福音的人，当他说\Quote{我不相信}时，你们如何回答他？至于我，除非因为天主教会权威的推动，我不会相信福音。所以，当我因其权威而同意相信福音的那些人告诉我不相信摩尼时，我怎能不同意呢？你们选择吧。如果你们说：相信天主教徒；他们的建议是不要信任你们；因此，相信他们，我就被禁止相信你们——如果你们说：不要相信天主教徒；你们就不能公平地使用福音来引导我相信摩尼；因为正是在天主教徒的命令下我才相信了福音——再者，如果你们说：你们当初相信夸奖福音的天主教徒是对的，但相信他们责难摩尼是错的：你们以为我如此愚蠢，以至于毫无理由地根据你们的喜好或厌恶而相信或不相信吗？因此，对我而言，远为公平和安全的是，既然已经在一次事例中相信了天主教徒，就不要投奔你们，直到你们不是命令我相信，而是以最清晰、最明了的方式使我理解某些事物为止。为了说服我，你们必须把福音放下。如果你们坚持福音，我就坚持那些命令我相信福音的人；并且听从他们，我根本不会相信你们。但如果你们碰巧在福音中找到一个无可辩驳的、为摩尼的宗徒身份作证的证据，你们就会削弱我对那些命令我不相信你们的天主教徒权威的尊重；其结果是，我将不能再相信福音，因为我是通过天主教徒才获得对福音的信仰；这样，你们从福音带来的任何东西都不再对我有分量。因此，如果在福音中找不到摩尼宗徒身份的明确证据，我将相信天主教徒而不是你们。但如果你们从那里读到一段显然有利于摩尼的经文，我将既不相信他们也不相信你们：不相信他们，因为他们关于你们对我撒了谎；也不相信你们，因为你们向我引用了那部我曾基于那些说谎者的权威而相信过的圣经。但我绝不会不相信福音；因为我相信它，却找不到也相信你们的方法。因为那里记载的宗徒名字，不包括摩尼的名字。而叛徒犹大的继任者是谁，我们在《宗徒大事录》中读到；\BibleRef{宗 1:26}如果我相信福音，就必须相信这书，因为这两部著作都是天主教的权威向我推荐的。同一卷书中载有关于保禄蒙召和宗徒职务的众所周知的故事。\EditorialNote{宗徒大事录第九章}现在，如果你们能，就请给我读一下福音中哪里说到摩尼被召为宗徒，或者任何其他我宣称相信的书。你们要读那段主应许将圣神作为护慰者赐予宗徒们的经文吗？关于那段经文，请看有多少重大事物阻止并威慑我相信摩尼。\par

% structure: p0013 source=h2 target=section reason=relative-heading-rank; base=h2

\section{第六章——摩尼为何自称基督的宗徒}

7. 我真是不明白，这封书信为何以\Quote{摩尼，耶稣基督的宗徒}开头，而不是\Quote{护慰者，耶稣基督的宗徒}。或者，如果基督所派遣的护慰者派遣了摩尼，那么为什么我们读到的是\Quote{摩尼，耶稣基督的宗徒}，而不是\Quote{摩尼，护慰者的宗徒}？如果你说基督自己就是圣神，那你就与圣经本身矛盾，因为主说：\BibleQuote{我必要求父，父就赐给你们另一位护慰者。}\BibleRef{若 14:16} 再者，如果你以基督的名称为正当，并非因为基督自己就是护慰者，而是因为他们本体相同——即不是因为他是一位，而是因为\OriginalTerm{是一体}{non quia unus est, sed quia unum sunt}——那么保禄也可以说\Quote{保禄，天主圣父的宗徒}；因为主说：\BibleQuote{我与父原是一体。}\BibleRef{若 10:30} 保禄从没有说过这样的话；也没有任何一位宗徒自称是圣父的宗徒。那么为何要有这种新花样？这难道不是某种欺骗的意味吗？因为他若认为没有区别，为何不在一部分书信中自称是基督的宗徒，在另一部分中自称是护慰者的宗徒呢？但在我所知的每一封书信中，他都写\Quote{基督的}；从未写过\Quote{护慰者的}。我们猜想这是为什么，岂非因为那骄傲——一切异端之母——驱使这人想要显得好像是被护慰者所派遣，却又与之联合得如此密切，以至自己取得了护慰者的名号？正如耶稣基督这人不是由天主子——即天主的德能和智慧——所派遣（万物藉他而造），而是根据公教信仰，被接纳到如此密切的关系中，以致他自己就是天主子——即在他内，天主的智慧在医治罪人中彰显出来——同样，摩尼希望人们认为，他是如此被基督所应许的圣神所接纳，以致我们从此就理解为，摩尼和圣神这两个名字同样意指耶稣基督的宗徒——即由耶稣基督所派遣的那一位，基督曾应许要派遣他。这是何等惊人的狂妄！真是无法言喻的亵渎！\par

% structure: p0015 source=h2 target=section reason=relative-heading-rank; base=h2

\section{第七章——摩尼的追随者以何种意义相信他是圣神}

8. 此外，你应当解释，尽管父、子、圣神在本质上是平等的（正如你也承认），你却不羞愧地谈论摩尼——一个被接纳与圣神联合的人——是由普通生育所生；然而你却不相信那个被接纳与独生智慧联合的人是由童贞女所生。如果人的肉体、\OriginalTerm{与男人同寝}{concubitus viri}、女人的子宫都不能玷污圣神，那么童贞女的子宫又怎能玷污天主的智慧呢？那么，这个摩尼，他自夸与圣神有关联，并在福音中被提及，他必须提出他对这两件事之一的要求——即他是被圣神所派遣，或是被接纳与圣神联合。若是被派遣，就让他称自己为护慰者的宗徒；若是被接纳联合，就让他承认那被独生子所接纳的人有一位人类的母亲，因为他承认在圣神所接纳的人身上同样有父亲和母亲。让他相信天主的圣言并未被玛利亚的童贞子宫玷污，因为他劝我们相信圣神不能被其父母的婚姻生活玷污。但如果你说摩尼是与圣神联合，并非在子宫内或受孕前，而是出生之后，你仍必须承认他具有从男人和女人而来的血肉本性。既然你不怕谈论摩尼的血和肉体本质来自普通生育，或他肉体中所包含的内在污秽，并且认为那接纳了这个人的圣神（如你所相信的）没有被这一切玷污，那么我又何必害怕谈论童贞女的子宫和身体是纯洁的，反而宁愿相信那与这个人联合的天主智慧，在他母亲的肉体中依然保持无瑕无污？因此，无论你的摩尼声称是被护慰者派遣，还是与护慰者联合，这两个说法都不能成立，我便保持警惕，拒绝相信他的使命或他的接纳。\par

% structure: p0017 source=h2 target=section reason=relative-heading-rank; base=h2

\section{第八章——摩尼的诞辰节}

9. \Person{摩尼}加上\Quote{因圣父天主的眷顾}这句话，还有什么意图呢？他不就是想借用了\Person{耶稣基督}（他自称是其宗徒）和圣父天主（他说自己是由圣子奉圣父眷顾而派遣）的名号后，让我们相信他自己就是圣神，即第三位吗？他的原话是：\Quote{\Person{摩尼}，\Person{耶稣基督}的宗徒，因圣父天主的眷顾。}圣神并未被提及，而一个援引\OriginalTerm{护慰者}{Paraclete}的应许来支持自己宗徒身份的人，本应特别提到圣神，好让他能借福音的权威说服无知的人。对此，你们当然会说：在宗徒\Person{摩尼}的名号中，我们已经有了圣神、护慰者的名号，因为祂屈尊进入了\Person{摩尼}。所以我再问：你们为何要攻击天主教会的教义，即那位神圣智慧所在者是由童贞女所生，而你们却毫不顾忌地宣称那位你们说圣神降临于他的人是经由正常生育而出呢？我不得不怀疑，这个\Person{摩尼}利用\Person{基督}的名号来进入无知者的心灵，是希望他自己被当作\Person{基督}本人来崇拜。我将简要说明这一猜测的理由。在我还是你们学说的学生时，我经常问：为什么主的逾越节庆日普遍无人问津，虽然偶尔有少数懈怠的崇拜者，但守夜、特别的斋戒以及任何特殊的礼仪均未规定——简而言之，没有任何特别的仪式——而你们却对你们的\OriginalTerm{贝玛节}{Bema}——即\Person{摩尼}被杀的日子——致以极大的尊荣，你们会搭建一座带精美台阶的高台，铺上珍贵的布料，醒目地置于会众面前。得到的回答是：应当庆祝的是那真正受苦者的受难日，而\Person{基督}并未出生，只是以虚幻的肉身形象出现在人眼前，他仅仅假装受苦，并未真正承受苦难。这难道不可悲吗？那些希望被称为基督徒的人害怕童贞女的子宫会玷污真理，却不怕谎言？但回到正题：凡是留心的人，谁能不怀疑\Person{摩尼}否认\Person{基督}是由女人所生、具有人身的目的，是为了让\Person{基督}的受难（如今其时日已成为全世界的大庆节）不被\Person{摩尼}的信徒所庆祝，从而减损他期望引导对自身死亡之日的隆重纪念的热忱呢？因为对我们来说，贝玛节在逾越节期间举行具有极大的吸引力，我们曾更加热切地渴望那个庆日（贝玛节），因为另一个曾极为甘甜的庆日已被收回了。\par

% structure: p0019 source=h2 target=section reason=relative-heading-rank; base=h2

\section{第九章 圣神何时被派遣}

10. 也许你会对我说：主所预许的\OriginalTerm{护慰者}{Paráklētos}何时来临？关于此事，即使我对此没有其他可信之物，我宁愿期待尚未到来的护慰者，也不愿承认祂已临于\OriginalTerm{摩尼}{Manichaeus}身上。然而，既然圣神的降临在\WorkTitle{宗徒大事录}中叙述得极其清楚，我又何必如此无端地冒险去相信异端者呢？因为在\WorkTitle{宗徒大事录}中记载如下：\BibleQuote{德敖斐罗！我在第一部书中，已论及耶稣所行所教的一切，直到祂藉圣神选择了宗徒，并命令他们宣讲\WorkTitle{福音}的那一天。祂在受难以后，用许多凭据向他们显示自己还活着，四十天之久显现给他们，讲论天主国的事。祂与他们一起进食时，吩咐他们不要离开耶路撒冷，但要等候父的恩许，即你们听我所说过的：\InnerQuote{若翰固然以水施洗，但你们却要以圣神受洗，并且你们不久以后，即在五旬节，将要领受圣神。}他们聚集时，就问祂说：\InnerQuote{主，此时祢要复兴以色列国吗？}祂对他们说：\InnerQuote{父凭自己的权柄所定的时候和日期，不是你们所能知道的。但当圣神降临于你们身上时，你们将充满圣神的能力，并要在耶路撒冷及全犹太和撒玛黎雅，直到地极，为我作证人。}} \BibleRef{宗 1:1} 你看，主在此提醒祂的门徒们父的恩许，即他们曾从祂口中听到的圣神的来临。现在让我们看看祂何时被派遣；因为不久之后我们又读到如下记载：\BibleQuote{五旬节到了，门徒们都聚集在一处。忽然，从天上传来一阵响声，好像暴风刮来，充满了他们所在的整个房屋。有些散开好像火舌似的舌头，停留在他们每人头上，他们都充满了圣神，照圣神赐给他们的话，说起外方话来。那时，住在耶路撒冷的有从天下各国来的虔诚犹太人。这声音一响，许多人都聚集起来，都很惊愕，因为人人都听见门徒们用各人的本地语言说话。他们都惊讶奇怪，彼此说：\InnerQuote{这些说话的难道不都是加里肋亚人吗？怎么我们每人却听见他们用我们出生地的语言说话呢？我们中间有\Person{帕提亚人}、\Person{玛待人}、\Person{厄蓝人}、\Place{美索不达米亚}、\Place{亚美尼亚}、\Place{卡帕多细亚}、\Place{本都}、\Place{亚细亚}、\Place{夫黎基雅}、\Place{旁非利亚}、\Place{埃及}、\Place{基勒乃}附近的\Place{阿非利加}地区、以及\Place{罗马}的侨居者、\Person{犹太人}和归化者、\Person{克里特人}和\Person{阿拉伯人}，他们都听见他们用我们自己的语言讲论天主的奇事。}他们都惊讶不已，彼此猜疑，说：\InnerQuote{这事是什么意思？}另有些人却讥笑说：\InnerQuote{他们喝醉了新酒。}} \BibleRef{宗 2:1} 你看圣神何时来临。你还要什么？如果\WorkTitle{圣经}可信，难道我不应最乐意相信这部\WorkTitle{宗徒大事录}吗？它拥有最强大的见证，并且已经广为流传、传承并与包含圣神许诺的\WorkTitle{福音}本身一起公开教导，连我们也相信这部福音。因此，读了这部\WorkTitle{宗徒大事录}——它在权威上与\WorkTitle{福音}并列——我发现圣神不仅被预许给这些真宗徒，而且如此明显地降临，以至在这个问题上毫无错误的余地。\par

% structure: p0021 source=h2 target=section reason=relative-heading-rank; base=h2

\section{第十章——圣神两次赐予}

11. 因为我们的主在人间受到光荣，就是祂从死者中复活和升天。因为若望福音记载：\BibleQuote{圣神尚未赐下，因为耶稣尚未受到光荣。} \BibleRef{若 7:39} 如今，如果圣神未赐下的原因是耶稣尚未受到光荣，那么耶稣一受到光荣，圣神立即赐下。既然这光荣是双重的——就人而言和就天主而言——圣神也就两次赐下：一次是在祂从死者中复活后，向门徒们嘘气，说：\BibleQuote{你们领受圣神吧！} \BibleRef{若 20:22}；另一次是在祂升天十天之后。这数字十代表完美；因为在涵盖一切受造物的数字七上，加上了创造者的\OriginalTerm{圣三}{Trinitas}。关于这些，属神的人有许多虔敬而谨慎的论述。但我必须紧扣主题；因为我目前的任务不是教导你——你可能会认为这是傲慢——而是扮演一个探询者的角色，并向你学习，正如我九年来徒劳尝试的那样。因此，现在我有一份关于圣神降临的文件可以相信；如果你叫我不要相信这份文件，正如你一贯的建议是不要盲目轻信而不加考虑，那我就更不会相信你的文件了。那么，抛开一切书籍，用逻辑的清晰揭示真理，使我的心中不再有疑惑；或者提出一些书籍，我在其中将发现不是专横地要求我信服，而是我可以学习的可靠陈述。也许你会说，这封书信也属于此类。那么，让我不再停留在门槛上：让我们看看内容吧。\par

% structure: p0023 source=h2 target=section reason=relative-heading-rank; base=h2

\section{第十一章 \OriginalTerm{摩尼}{Manichaeus} 许诺真理，却不履行诺言}

12. \Quote{这些，}他说，\Quote{是来自永活长流的泉源的健康言语；凡听见它们，并先相信它们，然后遵守它们所阐述的真理的人，将永不受死亡，反而要在荣耀中享受永生。因为那接受了这天上的知识的人，实在可称为有福的，他借此得以自由，并要在永远的生命中安息。}正如你所见，这是对真理的承诺，而非真理的赐予。你们自己也能轻易看出，任何错误都可以用这种方式包装，在华丽外表的掩护下悄然潜入无知者的心中。倘若他说：\Quote{这些是来自毒泉的有害言语；凡听见它们，并先相信它们，然后遵守它们所阐述的，将永不得复生，反而要像罪犯一样遭受悲惨的死亡：因为那陷入这地狱般错误的人，确实应被宣告为可悲的，他将沉沦其中，永远受苦。}他若这样说，便说了真理；但这样非但不能为他的书赢得任何读者，反而会在一切得到此书的人心中激起最大的厌恶。让我们继续看下文；也不要被那些好人与坏人、有学与无学的人都能使用的言词所欺骗。接下来是什么呢？\par

13. \Quote{愿无形天主的平安，}他说，\Quote{与真理的知识，临于那些既相信又服从神圣诫命的圣洁可爱的弟兄们。}阿们，我们应说。因为这祈祷是极为可爱可赞的。只是我们必须记住，这些话既可为假教师所用，也可为好教师所用。因此，若他所说不过如此，那么人人都可安全地阅读并接受它。我也不反对接下来的话：\Quote{愿光明的右手保护你们，从一切仇敌的攻击和世界的罗网中解救你们。}事实上，在这书信的开端，在到达其主要议题之前，我找不出任何毛病。因为我不愿在小事上浪费时间。现在，我们来看这位作者对他将论述之事的明确陈述。\par

% structure: p0026 source=h2 target=section reason=relative-heading-rank; base=h2

\section{第十二章——摩尼的狂想。世界组成前的争斗}

14. \Quote{关于那件事，}他说，\Quote{亲爱的帕提库斯弟兄，就是你所告诉我的，说你想知道亚当和厄娃诞生的方式，究竟是由一句话产生，还是从物质生出，我将按合宜的方式回答你。因为在许多著作和记载中，我们发现有不同作者的不同断言和不同描述。而关于此事的真实情况，对所有民族，甚至对那些长久并频繁讨论此事者，都是未知的。因为，倘若他们已清楚知晓亚当和厄娃的生成，他们就不会仍受制于腐朽与死亡了。}这就是向我们承诺对此事的清晰知识，好叫我们不再受制于腐朽与死亡。若这还不够，请看下文：\Quote{必须，}他说，\Quote{先说许多预备的话，然后才能毫无疑义地发现这一奥秘。}这正是我所要求的：要有真理的如此证据，能使我对它的认识毫无疑义。即便这一承诺并非由该作者自己作出，我也理当要求并坚持这一点，以便面对如此好处——完全清晰确定的真理——时，任何反对都不能让我羞于从公教基督徒变成摩尼教徒。现在，我们且听他陈述什么。\par

15. \Quote{因此，}他说，\Quote{若你愿意，请先听世界组成之前发生的事，以及争斗如何展开，好使你能够区分光明的本性与黑暗的本性。}这就是该作者所作的全然虚假和不可置信的陈述。谁能相信在世界组成之前有任何争斗发生？即便它是可信的，我们现在所想要的也是知道某事，而非相信。因为说波斯人和西徐亚人很久以前彼此争战，这是一个可信的陈述；但我们读到或听到时相信它，却不能凭经验或理智的真理知道它。那么，既然我被告知要得到某种东西——不是必须相信权威，而是毫无歧义地理解——我就要拒绝任何此类的陈述；更何况那些不仅不确定，而且不可信的陈述。但倘若他有某些证据使这些事情清晰可理解呢？让我们尽可能以全部的耐心和宽容来听下文。\par

% structure: p0029 source=h2 target=section reason=relative-heading-rank; base=h2

\section{第十三章——两种对立的实体。光明王国。摩尼教导不确定性而非确定性}

16. \Quote{起初，}他说，\Quote{这两种实体是分开的。光明帝国由天主父掌管，祂在神圣的根源中永恒，在德能中恢弘，在本性中真实，常为祂自己的永恒而喜乐，在祂自身内拥有智慧与生命的官能，借此祂也包括了祂光明的十二个肢翼，这些都是祂国度的丰沛资源。在祂的每个肢翼中，也储存着无数难以言喻的宝贵财富。但父自己，在赞美中为首，在伟大中不可测度，已将幸福而荣耀的世界与自己结合，其数目和延续不可计量；这位神圣而荣耀的父和元首与它们同在，在祂辉煌的领域中没有任何贫穷或软弱。这些无与伦比的领域如此建立于光明与福乐的区域之上，以致无人能够移动或扰乱它们。}\par

17. 这一切的证明在哪里？摩尼从哪里学到这些？不要以护慰者之名来惊吓我。因为，首先，我前来不是要相信未知之事，而是要按照你们自己所嘱咐我的谨慎，获得毋庸置疑的真理的知识。因为你们知道你们如何苦毒地嘲弄那些不加思考就相信的人。更有甚者，这位作者，在此开始讲述极其可疑之事，之前不久却已承诺给予完全而确凿的知识。\par

% structure: p0032 source=h2 target=section reason=relative-heading-rank; base=h2

\section{第14章——摩尼允诺不可置疑之事的认知，却又要求对可疑之事持有信德}

其次，若信德是向我要求的，我更愿固守\WorkTitle{圣经}，圣经告诉我，圣神曾来并默感了宗徒，主曾应许派遣祂。因此你必须证明，或\Person{摩尼}所说的是真的，从而向我阐明我无法相信之事；或\Person{摩尼}就是圣神，从而引导我相信你无法阐明之事。因为我宣认天主教信德，并藉此期望获得确切的认知。既然你试图推翻我的信德，你必当向我提供确切的认知，若你能的话，好使你指证我未经考虑而采纳了现今的信仰。你提出两个不同的命题——一为你说那发言者是圣神，另一为你说那发言者所教导的显然是真实的。我本可要求两个命题都提供不可否认的证据。但我不贪婪，只须被说服一个。证明这人是圣神，那么即使不理解，我也将相信他所说的是真实的；或证明他所说的是真实的，那么即使没有证据，我也将相信他是圣神。有比这更公平或更仁慈的吗？但你不能证明这两个命题中的任何一个。你找不到比赞美自己的信德而嘲笑我的信德更好的方法。所以，在我依次赞美我的信仰而嘲笑你的信仰之后，你认为就我们的判断与行为而言，我们会得出什么结果？岂不就是与那些允诺不可置疑之事的认知、却又要求我们对可疑之事持有信德的人分道扬镳？而我们却要追随那些邀请我们始于相信我们尚不能完全领悟之事的人。如此，被这信德本身所强化，我们或能达到一种状态，藉着心灵的内心光照与确认，得知我们所相信的，这不再来自人，而是来自天主本身。\par

18. 正如我曾请这位作者向我证明这些事，我现在问他自己是从哪里学来的。若他回答说，这些事是由圣神启示给他的，他的心智蒙受神性光照，得以知道它们是确定而明显的，那么他自己便指出了认知与相信之间的区别。认知属于那位完全知道这些事已被证明的人；但在那些只听到他讲述这些事的人那里，并没有传达认知，而只要求相信的默许。无论谁轻率地顺从，就成为摩尼信徒，并非由于认知无可置疑的真理，而是由于相信可疑的言论。我们不经事青年时就是如此被欺骗的。因此，\Person{摩尼}本应说这些事对他而言是清楚证明的，但那些听到他讲述的人必须相信他而无证据，而不是允诺认知、或清楚证据、或免除一切不确定的解答。但他若这样说，谁不会回答他说：若我必须不知而信，我为何不相信那些因学者与无知者一致赞同而广为人知、且在各民族中由最权威的权威所确立的事呢？因害怕有人这样对他说，\Person{摩尼}先是允诺确定之事的认知，然后又要求我们相信可疑之事，从而迷惑不经事的人。然后，若有人请他证明这些事对他自己已被证明，他又失败了，并叫我们也要相信这一点。谁能容忍这样的欺骗与傲慢？\par

% structure: p0035 source=h2 target=section reason=relative-heading-rank; base=h2

\section{第15章——摩尼的学说不仅不确定，而且虚假。他荒谬地幻想一片黑暗的领域和种族与神圣区域及天主的实体毗邻。首先，错误在于给天主的本性设定边界和界限，仿佛天主是一种具有空间广延的物质实体}

19. 若我藉着天主和我们的主之恩助，将证明这位作者的陈述不仅是可疑的，而且是虚假的，那又如何呢？还有什么比这更不幸的事呢：那迷信不仅未能传授它允诺的认知与真理，反而教导直接违背认知与真理的事？这从以下内容将更清楚显明：\Quote{在一方向，在这光明神圣之地的边界上，有一片广阔深邃的黑暗之地，其中住着火体、毁灭的族群。这里有无限的黑暗，从同一源头以不可计量的丰富流出，连同其本有的产物。在此之外是混浊泥泞的水及其居民；在这些水之内是可怕狂暴的风及其君王与祖先。再就是一片毁灭的火域，有其首领与子民。同样，在这片火域之内是一个充满烟与阴霾的族群，其中住着那可怕的君王与所有族群的元首，周围有无数君王，他自身是他们所有人的心智与源头。这就是那祸害之地的五种本性。}\par

20. 将天主说成是气态的甚或以太的身体，这在一切具有清明心智、拥有一定辨识能力的人看来是荒谬的；他们能感知智慧与真理的本性并非在空间中延展或分散，而是伟大的，且无需物质尺寸而赋予伟大，不在任何方向或多或少受限制，而是处处与万有之父同等广阔，并非此处有一样、彼处有另一样，而是处处完美，处处临在。\par

% structure: p0038 source=h2 target=section reason=relative-heading-rank; base=h2

\section{第16章——灵魂虽可变化，却无形体。它整个存在于身体的每一部分。}

然而，何必谈论那超越灵魂一切能力的真理与智慧呢？须知灵魂本身的本性虽被知晓为可变的，却仍然没有任何物质性的空间延展。因为任何由某种粗重物质构成的东西，必然可分割为部分，一部分在此处，另一部分在彼处。例如，\Quote{手指小于整个手，一根手指小于两根；这根手指有一个位置，那根手指有另一个位置，其余部分也是如此。}同样，在湿气中，较少量占据较小空间，较大量占据较大空间；一部分在杯底，另一部分在杯口。在空气中，每一部分也各有其位置；这间屋子里的空气不可能在同一时刻与邻居屋子里的空气共存。甚至就光本身而言，一部分从一扇窗户透入，另一部分从另一扇窗户透入；更多的部分从较大的窗户透入，更少的部分从较小的窗户透入。事实上，任何物质性实体，无论是天上的还是地上的、空气中的还是潮湿的，都不可能部分小于整体，也不可能同时使一部分占据另一部分的位置；而是，一部分在此处，另一部分在彼处，其空间延展是一种具有明确界限和部分（或曰片段）的实体。反之，灵魂的本性，即使我们不考虑其感知真理的能力，而只考虑其赋予身体统一性以及在身体内感觉的较低级能力，它似乎也不具有任何物质性的空间延展。因为当它完全呈现在任何感觉中时，它就在身体的每个分离部分中完全呈现。在手指中并不比在手臂中更小，如同手指的体积小于手臂那样；而是在各处数量相同，因为整体无处不在。因为当手指被触摸时，整个精神都感受到，尽管感觉并非通过整个身体。精神的任何部分都不会对触摸无意识，这证明了整体的临在。然而，它并非如此地呈现于手指或感觉中，以至于它放弃了身体的其他部分，或将自身聚集到发生感觉的那一处。因为当它完全呈现于手指的感觉中时，如果另一部分（例如脚）被触摸，它也同样完全呈现于这个感觉中：因此，在同一时刻，它完全呈现于不同的地方，而没有离开一处以便处于另一处，也没有一部分在此处，另一部分在彼处；而是通过这种能力，表明自己在同一时刻完全呈现于分离的地方。既然它完全呈现于这些地方的感觉中，这就证明它不受空间条件的束缚。\par

% structure: p0040 source=h2 target=section reason=relative-heading-rank; base=h2

\section{第17章——记忆容纳最大空间的范围的概念}

再者，如果我们考虑精神记忆的能力——不是针对理智的对象，而是针对物质对象，正如我们也见野兽能记住的一样（因为牲畜在熟悉的地方无误地找到路，动物回到它们的食槽，狗认出主人的面容，而且睡觉时它们常常低吼或吠叫，若非它们的精神保留了先前所见或通过某种身体感官所感知的事物的影像，这是不可能的）——谁能正确地设想这些影像在哪里被容纳、在哪里被保存、在哪里被形成呢？的确，如果这些影像不大于我们身体的尺寸，那么可以说精神在包含自身的身体空间中塑造并保留它们。然而，身体占据小的物质空间时，精神却旋转着广大空间的影像——天与地——丝毫不缺空间，尽管它们成群地来去；因此显然，精神并不扩散于空间：因为它不是被最大空间的影像所容纳，反而容纳它们；然而，并非在任何物质容器中，而是通过一种神秘的能力或力量，借此它能使之增大或缩小，能将其收缩于狭窄界限内或无限扩展，能随意整理或打乱，能增多或减少至少数或一个。\par

% structure: p0042 source=h2 target=section reason=relative-heading-rank; base=h2

\section{第18章——理智判断事物的真理及其自身的行动}

那么，关于感知真理的能力——当这些因身体感官印象而获得形状的影像与真理相悖时，能够有力地抵抗它们——又该说什么呢？这种能力辨别出（举一个具体例子）真正的\Place{迦太基}与它自己随意变换的想象中的\Place{迦太基}之间的区别，并且轻松自如。它表明\Person{厄壁鸠鲁}的无数世界（他的幻想在其中无拘无束地漫游）是由于同样的想象能力；不赘述例子，我们也从同一源泉获得光明之地（其无边的范围）以及黑暗种族的五个洞穴及其中的居民——\Person{摩尼}的幻想竟敢在那里冒用真理之名。那么，辨别这些事情的能力是什么呢？显然，无论其范围多大，它大于这一切，并且是在没有任何此类物质影像的情况下被设想的。如果你能找到，就为这种能力找一个空间吧；赋予它一种物质延展；给它一个巨大的身体。确实，如果你思考得好，你做不到。因为对于一切物质性的东西，你的精神都会就其可分割性形成看法，而且你可以随意使这些东西的一部分更大、另一部分更小；而借以判断这些事物的那个东西，你认识到它超越它们，不是在地点的高处，而是在能力的尊严上。\par

% structure: p0044 source=h2 target=section reason=relative-heading-rank; base=h2

\section{第19章——如果精神没有物质延展，那么天主更不会有}

21. 因此，既然心智如此容易变化——或因众多不相衬的欲望，或因期望之物的丰裕或匮乏而引起的感觉变化，或因想象力的无尽游戏，或因遗忘与记起，或因学习与无知——既然心智，我说，因这些及类似的原因而频遭变化，且被觉察为没有任何局部或物质的广延，并拥有一种超越这些物质条件的行动活力，那么对于天主本身，祂在不受纷扰与变化方面超越所有理智受造物，按各人所应得而赐予各人，我们又当如何思想或结论呢？心智敢于表达祂，却比看见祂更容易；而看见越清晰，表达的能力越弱。然而，如果如摩尼教的寓言不断断言的那样，这位天主在某一个方向上广延受限，而在其他方向上无限，那么祂就可以被每一个人随心所欲地以或大或小的许多部分或分数来度量；例如，二尺的广延比十尺的广延少八部分。因为这是所有在空间中有广延的性质的属性，因此它们不能全部在同一地方。但即使对于心智，情况也并非如此；这种对心智的贬低和歪曲的观点，存在于那些不适合这种探究的人中间。\par

% structure: p0046 source=h2 target=section reason=relative-heading-rank; base=h2

\section{第二十章——驳斥两个领域的荒谬观念}

22. 但是，或许与其这样对属血肉的心智说话，我们不如下降到那些人的观点中去，这些人要么不敢，要么尚不适合从对物质之物的思考转向研究非物质的、属神的本性，因此不能反思自己的反思能力，从而看到它如何对物质广延下判断，而自身却不具有广延。那么，让我们下降到这些物质观念中，并问问：按照摩尼的说法，在光明与神圣的领域的那一方和那一侧，黑暗领域位于何处？因为他谈到一个方向和一侧，却没有说是哪一侧，是右侧还是左侧。无论如何，很清楚，说到一侧就意味着有另一侧。但如果有三个或更多的侧，那么图形要么在所有的方向上都被限定，要么如果在一个方向上无限延伸，那么它在有侧的方向上必然是有限的。那么，如果光明领域的一侧有黑暗种族，那么它的其他侧或各侧是由什么限定的呢？摩尼教徒对此不回答；但当被追问时，他们说，在其它侧，他们所谓的光明领域是无限的，即在整个无垠的空间中延伸。他们没有看到，这是最愚钝的理解也能明白的：那样的话就没有侧了？因为侧是它被限定的地方。那么，他说，虽然那里没有侧？但你所指的\Quote{一个方向或一侧}，必然意味着存在另一个方向或侧，或者其它的方向或侧。因为如果只有一侧，你应当说\Quote{在侧}，而不是在 \Emph{一侧}；正如就我们的身体而言，我们恰当地说\Quote{用一只眼}，因为还有另一只；或\Quote{在一侧胸脯}，因为还有另一侧。但如果我们说到一个东西\Quote{在一个鼻子上}或\Quote{在一个肚脐上}，那么有学问和没学问的人都会嘲笑我们，因为只有一个。但我不在措辞上坚持，因为你们可能用了一个来表示唯一的一个。\par

% structure: p0048 source=h2 target=section reason=relative-heading-rank; base=h2

\section{第二十一章——如果光明领域与黑暗领域相连，它必须是物质的。与光明领域相连的黑暗领域的形状}

那么，什么与你们称为光明与神圣的领域的那一侧接壤呢？你们回答：黑暗领域。那么你们是否允许后者是物质性的呢？当然你们必须如此，因为你们断言所有形体都起源于它。那么，你们这些迟钝、属血肉的人，怎么没有看到，除非两个领域都是物质的，它们的一侧才能彼此相连？你们怎么会如此心智盲目，竟然说只有黑暗领域是物质的，而所谓的光明领域是非物质的、属神的呢？我的好朋友，让我们暂且睁开眼睛，趁现在有人告诉我们，看看最为明显的事实：两个领域如果都是物质的，它们的一侧才能相连。\par

23. 或者，如果我们太迟钝愚蠢而看不到这一点，就让我们听听：黑暗领域是否也有一侧，而在其他方向上是无限的，如同光明领域一样？他们不持守这一点，因为害怕使它显得与天主相等。因此，他们使它在深度和长度上是无限的；但向上，在它之上，他们坚持有无限的空虚空间。并且，为了不让这个领域显得是一个等于光明领域一半的分数，他们还在两侧收窄它。如同举最简单的例子：一张饼被切成四块正方形，三白一黑；然后设想三个白色方块连为一体，并设想它们向上、向下以及向后在所有方向上是无限的：这就代表摩尼教的光明领域。再设想黑色方块向下和向后是无限的，但向上有无限的空虚：这就是他们的黑暗领域。但这些是秘密，他们只向非常急切和焦虑的询问者透露。\par

% structure: p0051 source=h2 target=section reason=relative-heading-rank; base=h2

\section{第二十二章——光明领域的形态是两者中较差的}

那么，如果这样，黑暗领域明显有两侧接触光明领域。如果它有两侧接触，那么它必然接触两侧。这就是关于之前告诉我们的在一侧的情况。\par

24. 光明界的这个景象何其不雅！——像一座裂开的拱门，下方插入一个黑色楔子，仅在裂口方向有边界，其间有一虚空空间，无边空虚在黑暗界之上延伸。确实，黑暗界的形状比光明界更佳：因为前者裂开，后者被裂开；前者填补了后者中造成的缺口；前者其中没有虚空，而后者在所有方向都不确定，除了被黑暗楔子填满之处。出于一种无知且贪婪的观念，认为多个部分比单一整体更光荣，以致光明界应有六个部分，三个向上、三个向下，他们使这个界分裂开来，而不是分裂另一个界。因为，按此图形，虽可能没有黑暗与光的混合，却肯定有渗透。\par

% structure: p0054 source=h2 target=section reason=relative-heading-rank; base=h2

\section{第23章——拟人派不比摩尼派更糟}

25. 现在，比较一下——不是天主教会中属神的人，他们的心智在此生中尽可能领悟到神圣实体与本性没有物质延伸，也没有线条界定的形状；而是我们信仰中属血肉的软弱之人，当他们听到身体肢体被比喻地使用时，例如当说到天主的眼睛或耳朵时，他们习惯在幻想的自由中把天主想象为人形；把这些人与摩尼派比较，摩尼派惯于把他们的愚蠢故事当作伟大奥秘讲给焦虑的探询者听：然后判断谁对天主有更可容许、更可敬的观念——那些认为祂有最卓越人形的人，还是那些认为祂有无限物质延伸、但不是在所有方向，而是一方无限坚实、另一方裂开、有一空虚虚空、上方有不确定空间、黑暗界如楔子从下方插入的人？或者也许正确的说法是：祂在上方因自己的本性而不受限制，在下方却被敌对本性侵蚀。我与你一起嘲笑那些属血肉之人的愚昧，他们尚未形成属神的观念，以为天主有人形。你若能，也与我一起嘲笑那些不幸观念的人，他们把天主描绘成以如此不雅、不体面的方式裂开或切开，上方有那样的空隙，下方有那样可耻的削减。此外，还有这个区别：这些属血肉之人，若他们满足于从天主教会的乳房吮吸乳汁，不贸然投入草率意见，而是在教会中培养虔诚的探询习惯，在那里祈求以便领受，叩门以便为他们敞开，就会开始属神地理解圣经的比喻和寓言，逐渐察觉天主的能量适当地以耳朵、眼睛、手或脚，甚至翅膀和羽毛，盾牌、剑、头盔，以及所有其他无数事物的名称来表现。他们在这理解中进步越大，就越被确认是天主教徒。相反，摩尼派当他们抛弃物质幻想时，就不再是摩尼派。因为这是他们赞美摩尼的主要特别之点，即古代书籍中以比喻方式教导的神圣奥秘是为最后要来的摩尼保留，由他解决并证明；因此在他之后不会再有来自天主的教师，因为他没有以比喻或寓言说过任何东西，而是解释了那种古代说法，且他自己以简明直白的术语教导。因此，当摩尼派听到其创始人的这些话，在光辉神圣界的一边和边界上是黑暗界，他们没有解释可依赖。无论他们转向何处，他们自己幻想的悲惨束缚将他们带到最不雅的裂缝或突然停顿、连接或分裂，相信任何非物质本性（即使可变，如心灵）是真的也是骇人听闻的，更不用说天主的不可变本性了。然而，若我无法提升到更高事物，无法将我的思想从通过身体感官印在记忆上的虚假想象的缠绕中带出，进入精神存在的自由与纯洁，那么，认为天主有人的形状，岂不比把那个黑暗楔子钉在祂的下缘，并且因为无法覆盖上方无边的空虚而让它在无限空间中空荡闲置要好得多！有什么观念比这更糟？有什么更黑暗的错误能被教导或想象？\par

% structure: p0056 source=h2 target=section reason=relative-heading-rank; base=h2

\section{第24章——论摩尼派虚构中的本性数目}

26. 再者，我想知道，当我读到天主圣父和祂建立在光辉而幸福的领域上的王国时，圣父、祂的王国以及这领域，是否都具有同一本性与实体。如果是，那么分裂并穿透这领域的黑暗种族的楔子——这本身是一件难以言表的可憎之事——所施加的便不是天主的另一种本性或身体，而实际上是天主的本性本身遭受此等事。我恳求你们想想这一点：你们是人，请想想，并远离它；如果撕裂你们的胸膛就能从信仰中连根拔除这些亵渎的幻想，我祈求你们这样做。或者你们会说这三者并非同一本性，而是圣父是一种本性，王国是另一种，领域又是一种，各自有独特本性与实体，并按卓越程度排列？若果真如此，\Person{摩尼}本应教导四种本性，而非两种；或者若圣父与王国具有同一本性，而领域有其自身本性，他则应教导三种。若他只教导两种，因为黑暗领域不属于天主，那么光明领域在何种意义上属于天主？因为若它有其自身本性，且天主既未生它亦未造它，它便不属于祂，祂的王国的座落于属于他人之物上。或者若因其邻近而属于祂，则黑暗领域也必如此；因为它不仅与光明领域接壤，而且穿透它以至于将其一分为二。再者，若天主生它，它就不能有分离的本性。因为凡天主所生的，必与天主相同，正如天主教会对独生子所信的。如此，你必然回到那种骇人而可恶的亵渎：黑暗的楔子分裂的不是一个与天主有别而分离的领域，而是天主自身的本性。或者，若天主不是生，而是造它，祂是用什么造的呢？若是从自己，这与生有何区别？若是从另一种本性，这本性是善还是恶？若是善，则必有一种不属于天主的善本性，这你们几乎不敢断言。若是恶，则黑暗种族便不是唯一的恶本性。或者，天主是否取了那领域的一部分，将其转变为光明领域，以便在其上建立祂的王国？若是这样，祂会取全部，就不会有恶本性留存。因此，若天主不是用与自身不同的实体造了光明领域，祂必是从无中造它。\par

% structure: p0058 source=h2 target=section reason=relative-heading-rank; base=h2

\section{第二十五章——全能者创造不同程度的美善之物。关于两领域结合的任何描述，皆有不当或荒谬。}

27. 那么，如果你现在确信天主能够从无中创造某些善物，就请进入天主教会，并学习：所有天主按其卓越秩序从至高至低所创造并建立的本性，都是善的，且有些优于另一些；它们是从无中受造的，虽然天主——它们的造主——使用祂自己的智慧作为工具，可以说是，使无变为有，并且就它有而言，它是善的，而它的存在的有限性表明它不是由天主所生，而是从无中受造。如果你仔细考量，你会发现没有任何东西阻止你同意这一点。因为你不能使你的光明领域等同于天主，而不使黑暗的部分侵犯天主的本性本身。你也不能说它是天主所生，而不陷入同样的谬误，因必然结论是：作为天主所生，它必与天主相同。你也不能说它与天主有别，以免被迫承认天主将祂的王国置于不属于祂之物中，并有三种本性。你也不能说天主用与自身不同的实体造了它，而不在善之外又造出某种善，或在黑暗种族之外又造出某种恶。因此，你必须承认天主从无中造了光明领域；而你不愿相信这一点，因为如果天主能从无中造出某种伟大但次于祂自身的善，那么，既然祂是善的，且不吝惜任何善，祂也能造出次于前者的另一种善，再其次，依此类推，直到受造本性的最低之善，从而使整个总体——而非无限无量、无数量无尺度——具有固定而确定的一致性。再者，若你连这一点也不允许——即天主从无中造了光明领域——你将无法逃脱你的观点所导致的亵渎。\par

28. 或许，由于血肉的想象可以幻想任何它喜欢的形状，你或许能够为两领域的结合设计出某种其他形式，而不是向心灵呈现如此不悦且痛苦的描述——即天主的领域（无论是否与天主同一本性，至少天主的王国建立在那里）浩瀚无垠地铺展开来，形成一个巨大的团块，其部分松散地延伸至无限，而在其较低部分，那个无限大小的黑暗领域楔子侵入了它们。但无论你为两领域的结合设计出何种其他形式，你都无法抹去\Person{摩尼}所写的。我指的不是其他更详细描述的论著——或许因为它们仅在少数人手中，可能不那么困难——而是我们现在正考虑的这部\WorkTitle{基本书信}，你们所有被称为\Quote{光明者}的人通常都很熟悉它。这里的文字是：\Quote{在光辉神圣的领域边界的一侧是黑暗领域，深邃且无边无际。}\par

% structure: p0061 source=h2 target=section reason=relative-heading-rank; base=h2

\section{第二十六章——摩尼教徒被迫选择一种曲折、弯曲或笔直的结合线。第三种线将赋予适合两领域的对称与美。}

还有什么更多益处？我们已经听到了边界的情况。随你做出什么形状，画出何种线条，可以肯定的是，这无边无际的黑暗领域与光明领域的接合，必须是通过直线、曲线或弯曲的线之一。如果接合线是弯曲的，光明领域的边也必须弯曲；否则，其直边与弯曲边接合就会留下无限深的间隙，而非如先前所告诉我们那样，在黑暗之地之上只有虚空。如果存在这样的间隙，那么对于光明领域而言，最好让它更远离一些，在中间有更大的虚空，这样黑暗领域就完全接触不到它！那时，可能有一个无底深谷，以至于当那个族群中发生任何祸害时，尽管黑暗首领可能有鲁莽的愿望要跨过去，他们也会一头栽进深谷（\Quote{因为身体没有空气承载就不能飞行}）；而且由于下方是无限的空间，他们不能再造成伤害，尽管他们可能永远活着，因为他们将永远坠落。再者，如果接合线是曲线，那么光明领域也必须有一个相应的曲线的畸变。或者，如果黑暗之地向内弯曲得像一个剧场，那么光明领域中的相应线条也会有同样多的畸变。或者，如果黑暗领域有一条曲线，而光明领域有一条直线，它们不可能在所有点上接触。当然，如我之前所说，它们不接触更好，如果它们之间有这样一个间隙，使得各领域能保持截然分离，鲁莽的作恶者会一头栽下，从而无害。那么，如果接合线是直线，当然就不会再有间隙或沟槽，反而是一种如此完美的接合，使得两个领域之间达到最大可能的和平与和谐。有什么比一边以直线与另一边相遇，没有弯曲或断裂来扰乱贯穿无尽空间与无尽时间的自然且永恒的联系，更为美丽或更为适宜呢？即使存在分离，两端的直边本身也是美的，因为它们是直的；此外，即使有一段间距，它们作为平行线，尽管不相交，也会赋予两者对称性。加入接合后，两个领域变得完全规则与和谐；因为在描述或想象中，没有什么比两条直线的接合更美了。\par

% structure: p0063 source=h2 target=section reason=relative-heading-rank; base=h2

\section{第二十七章。—— 从黑暗领域的直线之美可取其美而不损其实体；因此，恶既不加也不减灵魂的实体；其边的直是造物主天主赐予黑暗领域的一种善。}

对于不幸的心智该怎么办呢？它们在错误中执迷不悟，被习惯紧紧束缚。这些人不知道自己在说什么，因为他们不思考。听我说：没有人强迫你们，没有人跟你们争吵，没有人用过去的错误嘲讽你们，除非某人没有经历过从错误中得救的天主慈悲：我们所渴望的，只是错误在某个时候被摒弃。不带仇视或苦毒，稍微想一想。我们都是人：让我们憎恨的不是彼此，而是错误和谎言。稍微想一想，我恳求你们。仁慈的天主，求祢帮助他们思考，在寻求者心中点燃真光。如果有什么是明显的，难道不是：正义比错误更好吗？那么，给我一个平静而安宁的回答：将连接于光明领域直线的黑暗领域直线弄弯曲，难道不会减损其美吗？如果你们不是固执不化，你们必须承认，不仅因为它被弄弯曲而失去了美，而且还失去了它本可以从与光明领域直线连接而获得的美。那么，在这种美的损失中——其中正直变弯曲，和谐变不谐，一致变不一致——是否损失了实体呢？那么，由此得知，实体并非恶；就如身体中，通过变坏的形式，美失落了，或者说减少了，以前被称为美丽的说成丑陋，令人愉悦的变成令人不悦；同样，在心灵中，正直意志的优美——它成就正义与虔敬的生活——当意志变坏时就受到损伤；通过这种罪，心灵变得悲惨，而不是如以前那样享受来自正直意志装饰的幸福，而实体没有任何增益或损失。\par

再想一想，尽管我们承认黑暗领域的边界因其他原因（如它是昏暗的、黑暗的）是恶的，但它在直这一方面并不是恶的。那么，如果我承认它的颜色中有某种恶，你们也必须承认它的直中有某种善。无论这善有多少，我们都不能将其归于任何其他，只能归于造物主天主，我们必须相信，一切自然中的善都来自祂，若我们要逃避致命的错误。那么，说这个领域是纯粹的恶，而它的直线边界中却发现了不少物质美的善，这是荒谬的；并且，将这个领域说成与全能而良善的天主完全隔绝，而我们在其中发现的这种善除了归于一切善的创造者以外不能归于任何其他，这也是荒谬的。但是，我们也被告知，这个边界是恶的。好吧，就算它是恶：如果它弯曲而不是直，它肯定会更坏。那怎么能是恶的极致，而比它自己更坏的东西可以被想象呢？变得更坏意味着有某种善，其缺乏使事物更坏。在这里，缺乏直会使线更坏。因此，它的直是某种善。你永远无法回答这个善从何而来的问题，除非你归诸祂，我们必须承认一切善，无论大小，都来自祂。但现在我们要从考虑这个边界转向其他事了。\par

% structure: p0066 source=h2 target=section reason=relative-heading-rank; base=h2

\section{第二十八章——\Person{摩尼}将五种本性置于黑暗区域}

31. \Quote{在那区域有火性的身体，毁灭性的种族，}他说。通过说到\InnerQuote{居住}，他必定是指那些身体是有生命且活着的。但是，为了不显得在吹毛求疵，让我们看看他如何将这区域的所有居民分成五类。\Quote{这里，}他说，\Quote{是无边的黑暗，从同一源头流出，不可计量的丰富，带着它所应产生的产物。在这之外是泥泞混浊的水，连同其居民；在其内有可怕而猛烈的风，连同其王和其祖先。然后，又是一个毁灭性的火区域，连同其首领和人民。类似地，在这之内是一个充满烟和阴霾的种族，那里住着可怖的王和万物的首领，周围有无数首领，他本身是他们的思想和源头。这就是这有害区域的五种本性。}我们在这里发现五种本性被提及作为同一本性的一部分，他称之为有害区域。这些本性是黑暗、水、风、火、烟；他如此安排，使黑暗在最外面。在黑暗之内他放置水；在水之内是风；在风之内是火；在火之内是烟。而且这些每一种本性都有其独特种类的居民，同样是五种。因为对于问题是，是否所有都只有一种，还是不同的种类对应不同的本性；回答是，它们是不同的：正如在其他书中我们发现记载说，黑暗有蛇；水有游动的受造物，如鱼；风有飞行的受造物，如鸟；火有四足动物，如马、狮子等；烟有两足动物，如人。\par

% structure: p0068 source=h2 target=section reason=relative-heading-rank; base=h2

\section{第二十九章——对这番荒谬之论的反驳}

32. 那么，这是谁的安排呢？谁作了区分和分类？谁赋予了数目、性质、形式、生命？因为这些事物本身都是善的，每种本性若非从天主——一切美善之物的创造者——获得，就不可能拥有它们。因为这不像诗人们关于想象的混沌的描述或假想，如同一无形之体，毫无形式、性质、度量、重量和数目、秩序和多样化；一种混乱的东西，完全缺乏性质，所以有些希腊作家称之为\OriginalTerm{无性质}{\GreekText{ἄποιον}}。\Person{摩尼}对他们所谓的黑暗区域的描述远非如此，以至于以完全相反的方式，他们将边与边相连，界与界相接；他们数出五种本性；他们分别、安排并赋予每种本性其固有的性质。他们也不让这些本性荒芜或废弃，而是用合适的居民充满它们；并且对这些居民，他们给予适当的形式，适应其居住之所，此外还赋予一切恩赐中首要的——生命。列举如此美好之物，却说它们与天主——一切美善之物的创造者——毫无关系，就是忽略事物中秩序的卓越，以及导至这一结论的错误的巨大危害。\par

% structure: p0070 source=h2 target=section reason=relative-heading-rank; base=h2

\section{第三十章——\Person{摩尼}置于黑暗区域的那些本性中美好之物的数目}

33. \Quote{但是，}有人回答道，\Quote{那五种本性的居住者之类别是凶暴和毁灭性的。} 好像我在称赞它们的凶暴和毁灭性似的。你看，我与你一同谴责你归给它们的恶事；请你与我一同称赞你归给它们的善事：这样就会显明，在你所谓穷凶极恶的境域中，善与恶是混合的。如果我与你一同谴责这区域中有害的，你必要与我一同称赞有益的。因为没有某种有益的影响，这些存在物就不能产生、被养育或继续居住在那区域。我与你一同谴责黑暗；请你与我一同称赞生产力。因为尽管你称黑暗为不可测量的，你却提到\Quote{适当的产物}。黑暗确实不是实在的实体，它无非是光明的缺席，正如赤身意味着缺少衣服，空虚意味着缺少物质内容：因此黑暗不能产生任何东西，尽管一个在黑暗中——即在光明缺席中——的区域可能产生一些东西。但暂且不论这一点，可以肯定的是，凡有产物出现的地方，必然有物质之有益的适应，以及受造存在物的各个部分在统一中的对称排列与构建——一种使它们彼此协调的智慧调整。谁能否认，这一切比黑暗更应受称赞，而非黑暗更应受谴责呢？如果我与你一同谴责水的浑浊，你必要与我一同称赞水，只要它们拥有水的形式和性质，以及水中游泳的居住者各部分的一致，它们维持和支配身体的生命，以及为健康之益处的每一种物质的特殊适应。因为尽管你指责水是浑浊和泥泞的，但你在承认它们具有产生和维持其活的居住者的性质时，暗示了存在着某种身体形状和部分的相似性，赋予统一和谐的特性；否则根本不可能有身体：而且，作为一个理性存在者，你必须看到这一切都是值得称赞的。并且，无论你将这些居住者的凶残程度，及其攻击中的屠杀和毁灭夸大到何等地步，你仍然为它们保留着形式的规律界限，藉此每个身体的各部分得以相互协调，以及它们有益的适应，以及生命原则的调节力量，将身体的各部分友善和谐地结合在一起。并且，如果用常识来看待这一切，就会发现它们更值得称赞，而非缺点更应受谴责。我与你一同谴责风的可怕；请你与我一同称赞它们的本性，因为它们提供呼吸和滋养，以及它们的物质形式在其各部分连接中的连续性和扩散：因为这些，这些风拥有产生、滋养并保持你所提到的这些居住者的活力的能力；并且在这些居住者中——除了在所有有生命的受造物中已经称赞的其他事物之外——还有这种特殊的能力，能迅速轻易地来去自如，以及飞行时翅膀的和谐拍动和规律运动。我与你一同谴责火的毁灭性；请你与我一同称赞这火的生产力，以及这些产物的生长，以及火对所产生的存在物的适应，使它们具有凝聚力，在度量和形状上达到完美，并且能够生活并居住在那里：因为你看，这一切都值得惊叹和赞美，不仅在如此宜居的火中，也在居住者身上。我与你一同谴责烟的稠密，以及正如你所说，居于其中的君王的野蛮性情；请你与我一同称赞这烟本身的所有部分的相似性，藉此它按照其本性保持其各部分之间的和谐与比例，并具有使其成为其自身的统一：因为没有人能够冷静思考这些而不会感到惊叹和赞美。此外，你甚至赋予烟生产和能量的能力，因为你说君王住在其中；因此在那个区域，烟是能生产的，这在这里从未发生，而且它为居住者提供了健康的居所。\par

% structure: p0072 source=h2 target=section reason=relative-heading-rank; base=h2

\section{第三十一章 同一主题续}

34. 即使在烟雾王子本身，你也不应只提到他的凶残这一恶质，难道不应该注意他本性中其他你不得不承认是可称赞的事物吗？因为他有灵魂和身体：灵魂赋予生命，身体被赋予生命。既然灵魂主宰，身体服从；灵魂引导，身体跟随；灵魂赋予稳定，身体不瓦解；灵魂赋予和谐运动，身体由匀称的肢体框架构成。在这唯一王子身上，你不是被迫承认那有序的和平或和平的秩序吗？适用于一人的也适用于其余所有人。你说他对他人凶残暴虐。这不是我所称赞的，而是你不愿注意的其他重要事物。这些事物，当被察觉和深思后——在经过任何轻率信任摩尼之人的劝告之后——会引领他清楚确信：在谈论那些本性时，他所说的在某种意义上都是善的事物，并非完美和不受造的，如同天主独一圣三，也非受造物中的较高等级，如同圣天使和永远蒙福的权能；而是最低的等级，按其禀赋的微小尺度排列。这些事物被无学问者与更高事物比较时，被认为应受责备；并且由于它们缺乏某种善，它们所拥有的善就被称为恶，因为它是残缺的。我这样讨论摩尼所列举的本性的另一个理由是，这些被命名的事物是我们在这世界上所熟悉的。我们熟悉黑暗、水、风、火、烟雾；我们也熟悉动物，爬行的、游水的、飞行的；四足的和两足的。除了黑暗（正如我已经说过的，黑暗只是光的缺乏，对它的感知只是视觉的缺乏，如同对寂静的感知是听觉的缺乏；并非黑暗是什么，而是光不存在，也并非寂静是什么，而是声音不存在），所有其他事物都是自然的性质，为众人所熟悉；那些本性的形式，就其存在而言是可称赞和善的，没有一个明智之人归之于除天主——一切善物的作者——以外的任何作者。\par

% structure: p0074 source=h2 target=section reason=relative-heading-rank; base=h2

\section{第32章——摩尼从可见物中获得其虚构观念的排列}

35. \Person{玛尼}在将那些从可见事物学来的本性，按其臆想安排以代表黑暗种族时，显然错了。首先，他使黑暗成为能生殖的，这是不可能的。但他回答说：\Quote{这黑暗不同于你所熟悉的那种。}那么，你如何能让我明白它呢？在许诺了这么多知识之后，你竟要我相信你的话吗？假设我相信你，至少这一点是确定的：如果黑暗没有形状（如通常的黑暗），它就不能产生任何东西；如果它有形状，那它就比普通黑暗更好。然而，当你称它不同于普通黑暗时，你愿意我们相信它更坏。你莫如说寂静（它之于耳如黑暗之于眼）在那个区域产生了一些聋哑的动物；然后，针对寂静不是一种本性的反驳，你又说那是不同于普通寂静的寂静。总之，你可以对被你误导相信你的人随意说什么。无疑，与动物生命起源相关的明显事实使\Person{玛尼}说蛇是在黑暗中产生的。然而，有些蛇视力锐利，喜爱光明，它们似乎提供了最有力的证据反对这种看法。然后，在水中游泳的生物很容易从这里获得想法，并应用于那个区域的臆想之物；同样，风中的飞行物也是如此，因为此世下层空气（鸟在其中飞翔）的流动被称为风。至于他关于火中四足动物的想法从何而来，无人知晓。尽管如此，他仍故意这么说，虽未经深思，且出于极大误解。通常给出的理由是，四足动物贪食而好淫。但许多人在贪食上超过任何四足动物，即使他们是两足动物，被称为\Quote{烟之子}，而非\Quote{火之子}。鹅也贪食如任何动物；尽管他可将之放入火中作为两足动物，或放入水中因其爱游泳，或放入风中因其有翅膀且有时飞翔，它们在这个分类中确实与火无关。至于好淫，我想他想到的是嘶鸣的马，它们有时咬断缰绳冲向母马；匆忙写作时，他心中想着这事，竟忘了普通的麻雀，与之相比，最烈的种马也显得冷淡。他们对两足动物归于烟所给的理由是，两足动物自负而骄傲，因为人出自此类；这个想法似乎合理，即烟如骄傲之人那样升入空中，圆而膨胀。这个想法可用于比喻骄傲者的描述，或寓言式的表达或解释，但不能使人相信两足动物生于烟中且由烟所生。他们同样有理由说人出生于尘土，因尘土常以相似的旋转和崇高运动升天；或生于云中，因云常以这样的方式从大地升起，使远处观看的人不确定它们是云还是烟。再者，在水和风的情况下，他为何使居民适合地方的特性，如我们在水中看到游泳之物，在风中看到飞行之物；而面对火和烟时，这个大胆的撒谎者却不羞愧地指派最不可能的居民到这些地方？因为火焚烧并消灭四足动物，烟令两足动物窒息而死。至少他必须承认，他在黑暗种族中使这些本性变得比这里更好，尽管他希望我们相信一切更坏。因为照此，那里的火产生并养育四足动物，给它们不仅无害而且最舒适的住所。烟也为自己慈爱的胸怀所生的后代提供空间，并养育它们到王子的等级。因此我们看到，这些谎言（它们增加了异端者的数目）起源于对世上可见事物的肉体感觉，只是缺乏谨慎和辨识力，感觉形成概念后，由想象孕育，然后被放肆地写下来并发表。\par

% structure: p0076 source=h2 target=section reason=relative-heading-rank; base=h2

\section{第三十三章——每一本性，作为本性，都是善的}

36. 但我们最想强调的论点是天主教教义的真理，即天主是一切性体的作者——如果他们能理解的话。我之前曾这样强调：\Quote{我与你一同谴责破坏性、盲目、浓密的混沌、可怕的暴力、朽坏性、王侯的残暴等；请与我一同称赞形式、分类、安排、和谐、结构的统一、对称与各部分的对应、对生命气息和滋养的供应、有益的适应、由心智进行的规整和控制、身体的服从，以及性体中居住者与被居住者的同化与一致，以及其他一切同类事物。} 由此，若他们诚心思索，便会明白即使在那些他们认为恶独立且完满的领域，也意味着善与恶的混合：因为如果上述的恶被除去，善仍会存留，毫无减损对其的赞美；反之，如果善被除去，就没有任何性体留下。因此，凡能看见的人都会看出：每一个性体，就其作为性体而言，都是善的；因为在同一件事物中，我发现了可赞之处，而他却发现了可责之处：若除去善，则无性体存留；但若除去那些可厌之处，性体仍会完好无损。从水中除去稠浊和混浊，剩下的就是纯净清澈的水；除去其各部分的凝聚，水就不复存在。因此，如果恶被除去后，性体以更纯净的状态存留，而善被除去后性体则完全不复存在，那么必定是善构成了事物所在其中的性体，而恶并非性体，而是与性体相悖。从风中除去你所指责的可怕和过度的强力，你可以设想温柔和缓的风；除去其各部分的相似性——这赋予它们物质的连续性——以及物质存在所必需的一致性，就没有任何可设想的性体存留。逐一举例将过于冗长；但所有不带派性偏见思考这个问题的人都必须看到：在他们列举的性体清单中，那些可厌之物是附加于性体的；当它们被除去时，性体比从前更好。这表明，性体就其作为性体而言是善的；因为当你从性体中除去善而非恶时，就没有任何性体存留。注意，你们这些希望得出正确判断的人，听听关于残暴王侯本身所说的话。如果除去他的残暴，看看会有多少优异之物存留：他的物质结构、对称的肢体、统一的形态、各部分的稳定连续性、作为统治和赋予生命的心智的有序调整，以及作为服从和受生命的身体。除去这些以及我可能遗漏的其他事物，就不会有任何性体存留。\par

% structure: p0078 source=h2 target=section reason=relative-heading-rank; base=h2

\section{第三十四章——性体不能没有某些善。摩尼教徒耽于谈论恶}

37. 但或许你会说这些恶不能从性体中除去，因此必须被视为自然的。目前的问题不在于什么能被除去、什么不能；但当我们看到善可以在没有这些恶的情况下被设想，而没有这些善则任何性体都无法被设想时，这确实有助于清楚认识：这些性体就其作为性体而言是善的。我可以设想没有泥泞搅动的水；但没有各部分稳定的连续性，任何物质形式都无法以任何方式被思想或感觉。因此即使是这样的泥水，若无作为其物质存在条件的善，也不能存在。至于回答说这些恶不能从这类性体中除去，我反驳说，善也不能从它们中除去。那么，为什么你们应该因那些你们认为不能除去的恶而称这些事物为自然的恶，却拒绝因那些已被证明不能除去的善而称它们为自然的善呢？\par

38. 你接下来或许会问——正如你惯常当作最后手段那样——这些我表示也予以谴责的恶从何而来。如果你先告诉我那些你也不得不称许的善物从何而来——假如你不想全然不合情理——我也许会告诉你。但我何必问这个呢，既然我们双方都承认，一切不论怎样的善物，无论多么伟大，都来自唯一的天主，那至善者？因此，你们自己就必须反对\Person{摩尼}，他将我们刚才提及并公正称许的这一切重要的善物——每一种本性中各部分的连续与和谐，有生命之物的健康与活力，以及其他不胜枚举的东西——（安置在一个想象的\OriginalTerm{黑暗区域}{region of darkness}里，从而将它们与那位他所承认的万善之源的天主完全分离）。他忽略了那些善物，只留意令人不快的东西；如同一个人被狮子的吼叫惊吓，看见它拖曳并撕裂所捕获的牲畜或人的尸体，竟因幼稚的怯懦而恐惧过度，只看见狮子的残忍与凶猛，忽略或不顾所有其他品质，便谴责这动物的本性不仅是恶的，而且是巨大的恶，他的恐惧更增添了言辞的激烈。但倘若他看见一头驯服的狮子，其凶猛已被驯服，尤其如果他从未被狮子惊吓过，他便会在没有危险和恐惧的情况下，有闲暇观察并赞叹这动物的美丽。我对此的评语是与我们的主题紧密相连的：任何本性在某些情况下可能令人不快，以致引发对整个本性的憎恨；但显然，一头真实活着的野兽的形状，即使在森林中引起恐惧，也远远好于墙上画作所称赞的人工仿制品。因此，我们不可被\Person{摩尼}误导而陷入这种错误，也不可因他对某些事物的挑剔——他这样挑剔以致让我们完全否定它们——而妨碍我们观察本性的形式，因为我们无法证明它们应受全面否定。当我们的心灵如此镇定并准备好作出公正的判断时，我们就可以问：那些我表示予以谴责的恶从何而来？如果我们把所有恶都归在一个名称下，就更容易看清这一点。\par

% structure: p0081 source=h2 target=section reason=relative-heading-rank; base=h2

\section{第三十五章——唯独腐败是恶。腐败不是本性，而是违反本性。腐败暗示先前的善}

39. 因为谁能怀疑所谓的恶无非就是腐败呢？的确，不同的恶可以用不同的名称来称呼；但在一切可感知恶的事物中，那恶就是腐败。所以，有教养心灵的腐败是无知；审慎心灵的腐败是不审慎；正义心灵的腐败是不义；勇敢心灵的腐败是怯懦；平静安宁心灵的腐败是贪欲、恐惧、忧伤、骄傲。同样，在有生命的身体中，健康的腐败是痛苦和疾病；力量的腐败是疲惫；休息的腐败是劳苦。再次，在任何有形物中，美的腐败是丑；直的腐败是曲；秩序的腐败是混乱；完整的腐败是分离、破裂或缩小。要一一列举这里所提到的事物以及无数其他事物的各种腐败名称，将是冗长而费力的；因为在许多情况下，这些词既可适用于心灵也可适用于身体，而且在无数情况下，腐败有它自己独特的名称。但已经说得足够多了，足以表明腐败只有通过取代自然状态才会造成损害；因此，腐败不是本性，而是违反本性。既然腐败是唯一能找到的恶，并且腐败不是本性，那么就没有本性是恶的。\par

40. 但如果，或许，你无法理解这一点，请再想想：凡是被腐败的，都被剥夺了某种善；因为如果它没有被腐败，它就是未腐；如果它无论如何不能被腐败，它就是不朽坏的。现在，如果腐败是一种恶，那么未腐和不朽坏都必须是善物。然而，我们目前并非在谈论不朽坏的本性，而是在谈论那些容许腐败的事物，它们虽然未被腐败，可以称为未腐，但并非不朽坏。唯独那不仅没有被腐败，而且任何部分都不能被腐败的，才能称为不朽坏的。因此，一切未腐但易于腐败的事物，一旦开始被腐败，就被剥夺了它们作为未腐时所拥有的善。这并非轻微的善，因为腐败是大恶。腐败的持续增加暗示着善的持续存在，它们可能被剥夺这些善。所以，假定存在于\OriginalTerm{黑暗区域}{region of darkness}的本性，要么是腐坏不了的，要么是易腐坏的。如果它们是不朽坏的，它们就拥有一个至高无上的善，没有比它更高的了。如果它们是易腐坏的，则要么被腐败了，要么没有被腐败。如果它们没有被腐败，它们就是未腐，说任何事物未腐就是给予它极大的赞誉。如果它们被腐败了，它们就被剥夺了未腐这一大善；但剥夺暗示着它们先前曾拥有被剥夺的善；而如果它们曾拥有这善，它们就不是恶的极致，因此所有摩尼的虚构故事都是虚假的。\par

% structure: p0084 source=h2 target=section reason=relative-heading-rank; base=h2

\section{第三十六章——恶的源头或善的腐败的源头}

41. 在这样探究了什么是\OriginalTerm{恶}{malum}，并得知它并非本性，而是反乎本性之后，我们接下来必须探究它从何而来。如果\Person{摩尼}这样做了，他本可免于陷入这些严重错误的陷阱。他不合时宜、不按顺序，从探究恶的起源开始，却没有先问恶是什么；因此他的探究只引导他接受了愚蠢的幻想，而深受肉体感觉滋养的心灵很难摆脱这些幻想。也许，有人不再渴望争论，而是渴望脱离错误，会问：我们发现在不朽坏的美善之物中常见的恶——\OriginalTerm{败坏}{corruptio}，从何而来？这样的人若以极大的热忱寻求真理，并以持续的勤勉虔诚地叩门，很快便会找到答案。因为虽然人可以用言语作为表达思想的某种符号，但教导乃是不朽坏的真理本身的工作，祂是唯一的真导师，唯一的内心导师。祂也成为了外在的导师，为要把我们从外在召回内在；祂取了奴仆的形体，为将祂的高度下降到那些向祂上升之人的认知中，祂屈尊以卑微向卑微者显现。让我们因祂的名祈求，并借着祂在做此探究时向圣父寻求仁慈。因为要用一句话回答\Quote{败坏从何而来？}这个问题，答案是：因为这些能够败坏的本性并非由天主所生，而是祂从无中创造的；既然我们已经证明这些本性是善的，没有人可以恰当地说它们被天主创造时并非善的。若有人说天主创造了它们成为完美美善，必须记住，唯一的完美美善就是天主自己，这些美善之物的创造者。\par

% structure: p0086 source=h2 target=section reason=relative-heading-rank; base=h2

\section{第37章 唯独天主是完美美善}

42. 你问，倘若那些事物也是完美美善的，会有何妨害？然而，若有人承认并相信天主圣父的完美美善，要探究我们应虔敬地将任何其他完美美善之物（假设其存在）归为何源，我们唯一正确的答复是：它来自天主圣父，祂是完美美善的。我们必须牢记，凡出自祂的，都是由祂所生，而非祂从无中创造，因此它像天主自己一样是完美美善的，即不朽坏的美善。所以我们看到，要求从无中创造的事物与那位由天主自己所生、并如同天主是独一的一样独一者同样完美美善，是不合理的；否则天主就会生出与祂不相似的东西。因此，为这位独生子寻找弟兄——万物都借着他由圣父从无中创造出来——是愚昧和悖逆的，除非是在这一点上：祂屈尊降生成人。因此，在圣经中，祂既被称为独生子，又被称为首生者；是圣父的独生子，又是从死者中首生者。\Person{若望}说：\BibleQuote{我们见过祂的光荣，正如圣父独生子的光荣，充满恩宠和真理。}\BibleRef{若 1:14} \Person{保禄}也说：\BibleQuote{使祂在众多弟兄中作首生者。}\BibleRef{罗 8:29}\par

43. 那么，你理性灵魂的本性啊，请服从，让自己略低于天主，但仅在此范围内低于，以致在祂之后，再无其他事物在你之上。我恳求你，服从并屈服于祂，免得祂将你降得更低，落入惩罚的深渊，在那里持续的惩罚将不断削减你所拥有的美善。你若因祂在你之先而愤怒，便是高举自己反对天主；你若因拥有这一美善——唯独祂在你之上——而未能以难以言表的喜乐欢欣，便是对祂存有极其悖逆的想法。这一点确定无疑之后，你不可说：\Quote{天主本应使我成为唯一本性：在我之后不应有任何美善之物。}在下一位于天主的美善之后，就不可能是最后一位。由此最清楚地看出，天主赐予了你何等大的尊荣：祂在本性的秩序中唯独统治你，却为你创造了其他美善之物供你管辖。不要惊奇它们现在并非事事服从你，有时甚至使你痛苦；因为你的主对于受你管辖的事物拥有比你自己更大的权柄，如同主人对于他仆人的仆人。那么，当你犯罪——即违背你的主——时，你先前所管辖的事物被用作惩罚你的工具，这有什么可奇怪的呢？因为有什么比天主更公义呢？这发生在人类的本性中，在\Person{亚当}身上，此处不宜详述。只要说，公义的统治者行事合乎祂的本性，无论是公正的赏报还是公正的惩罚，无论是生活正直者的幸福，还是加于罪人的刑罚。然而你并未被撇下没有仁慈，因为藉着预定的事物和时间的分配，你被召叫回头。如此，至高造物主的公义治理甚至延伸到尘世的美善之物上，它们被败坏又被修复，为使你得到与惩罚交织的安慰；为使你既因美善之物的秩序而喜悦时赞美天主，又因在邪恶的经验中受考验而投靠祂。因此，就尘世之物服从你而言，它们教导你你是它们的统治者；就它们使你痛苦而言，它们教导你要服从你的主。\par

% structure: p0089 source=h2 target=section reason=relative-heading-rank; base=h2

\section{第38章 天主创造的本性；败坏来自无}

44. 如此，虽然朽坏是一种恶，虽然它并非来自本性的造主，而是由于万物从虚无中被造，但在天主对其所造万物的治理与管束中，连朽坏也被如此安排：它只伤害最低下的本性，以惩罚被定罪者，并考验和教训回头者，使他们能亲近那不朽坏的天主，并保持不朽，这是我们唯一的善；正如先知所说：\Quote{但我亲近天主，于我有益。}你不可说：天主没有造可朽坏的本性；因为就其本性而言，是天主造了它们；但就其可朽坏而言，并非天主所造；因为朽坏不能来自那唯一不朽坏者。你若能领受这话，就感谢天主；若不能，就安静下来，不要谴责你尚未明了之事，而要谦卑地等候那心灵之光的主，好使你明白。因为在\Quote{可朽坏的本性}这一表达中有两个词，而非仅一个。同样，在\Quote{天主从虚无中造了}这一表达中，\Quote{天主}和\Quote{虚无}是两个分开的词。因此，将属于各自的归于各个词，使\Quote{本性}一词与\Quote{天主}相连，而\Quote{可朽坏}一词与\Quote{虚无}相连。然而，连朽坏本身，尽管其根源不在天主，仍须由祂按照无生命事物的秩序及祂有理智受造物的功过而管理。因此，我们正确地说，赏报和惩罚都来自天主。因为天主不造朽坏，这与祂将朽坏加给那配得朽坏的人，即那已开始犯罪而败坏自己的人，是一致的；这样，那故意屈服于朽坏诱惑的人，将违背自己意志而承受其痛苦。\par

% structure: p0091 source=h2 target=section reason=relative-heading-rank; base=h2

\section{第三十九章——恶以何种方式来自天主}

45. 不仅在旧约中写着：\Quote{我造善，也造恶；}更在新约中清楚地记载，主说：\Quote{不要怕那些杀身体而不能做更多的；但要怕那杀了身体以后，有权将灵魂投入地狱的。}而且，在神圣审判中，自愿的败坏附加了惩罚性的朽坏，这由宗徒保禄清楚宣告：\BibleQuote{天主的宫殿是圣的，这宫殿就是你们；谁败坏天主的宫殿，天主必要败坏那人。}\BibleRef{格前 3:17} 这话若是在旧律法中说出的，摩尼教徒会何等猛烈地抨击，说天主是败坏者！由于害怕这个词，许多拉丁译员译为\Quote{那人，天主必要毁灭}，以代替\Quote{败坏}，虽避开冒犯之词，但意义未变。尽管这些人会抨击旧律法或先知书中任何称天主为毁灭者的经文。但希腊原文显示\Quote{败坏}是正确词语；因为清楚写着：\Quote{谁败坏天主的宫殿，天主必要败坏那人。}若要求摩尼教徒解释这些话，他们为了免使天主成为败坏者，会说这里\Quote{败坏}的意思是\Quote{交付于朽坏}，或类似解释。假若他们本着这种精神阅读旧律法，便会在其中发现许多美妙之事；并且，对于他们不理解的章节，不是恶意攻击，而是虔敬地推迟探究。\par

% structure: p0093 source=h2 target=section reason=relative-heading-rank; base=h2

\section{第四十章——朽坏趋于非存在}

46. 但若有人不信朽坏来自虚无，就让他将存在与非存在摆在自己面前——仿佛各在一侧（如此说，是为不超越迟钝者的领会）；然后让他设定某物，比如动物的身躯，置于二者之间，并自问：当身体形成和产生、体积增大、获得营养、健康、力量、美丽、稳定时，就它的延续和持久而言，它是在趋向这一侧还是那一侧，是趋向存在还是非存在？他必不难看出，即便在初形中已存在，并且身体越在形态、形状和力量上建立和增长，就越趋于存在，趋向存在一侧。然后，再让身体开始朽坏；它的整个状态衰弱，活力萎靡，力量衰减，美丽消褪，骨骼离散，各部分的一致性瓦解崩坏；然后让他问：在朽坏中，身体现在趋向何处，是趋向存在还是非存在？他必不致如此盲目或愚蠢，以致怀疑如何回答自己，或看不出：事物越朽坏，就越接近消亡。但凡是趋于消亡的，就是趋于非存在。既然我们必须相信天主是永恒不变、不朽坏的存在，而所谓虚无显然完全是全然不存在的；既然在你将存在和非存在摆在面前之后，你已观察到：可见之物越增长，越趋向存在；越朽坏，越趋向非存在；那么，为何你面对任何一个本性时，却难于说明其中什么出自天主，什么出自虚无？因为可见的形态是本性的，而朽坏是反本性的。形态的增长导向存在，我们承认天主是至高存在；朽坏的增长导向非存在，我们知道那不存在的就是虚无。那么，我说，当你同时拥有\Emph{本性}和\Emph{可朽坏}这两个词时，为何你竟无法说明一个可朽坏的本性中，什么来自天主，什么来自虚无？为何你寻求一个与天主对立的本性？因为，你若承认祂是至高存在，那么非存在便是与祂对立的。\par

% structure: p0095 source=h2 target=section reason=relative-heading-rank; base=h2

\section{第四十一章——朽坏经天主允许，并来自我们}

47. 你问：为什么腐化从天主所赐予本性的东西中夺去？它什么也不夺取，除非是天主允许的；祂按公义且有秩序的审判，按照无知之物的层次和理智之物的功过而允许。说出的话作为感官对象消逝，沉寂于静默；然而这些消逝之言的来来去去构成了我们的言语，而有规律的静默间歇则带来悦耳且恰当的分辨；同样地，那些具有这种最低级美的暂时本性，正是经过变迁而存在，它们所产生之物的死亡赋予了它们个体性。如果我们的感官和记忆能正确领会这种美的秩序与比例，它就会如此令我们喜悦，以至于我们不敢把那些产生区别的不完美称为腐化。而且，当因损失所爱的暂时事物逝去，通过它们独特的美而使我们感到痛苦时，我们既为自己罪过付出代价，也被劝勉将情感寄托于永恒之物。\par

% structure: p0097 source=h2 target=section reason=relative-heading-rank; base=h2

\section{第四十二章 劝勉至善}

48. 那么，我们不要在这种美中寻求未给予它的东西（因缺乏我们所寻求的，这正是最低级的美）；而在它所被赋予的东西中，让我们赞美天主，因为祂甚至将这可见形式的大善赐予了最低级的美。并且，我们不要像情人那样依恋这种美，而要像赞美天主的人那样超拔于它；并从这个优越的位置对它作出判断，而不是如此束缚于它，以至与它一同被审判。让我们速速奔向那善，它没有空间的运动或时间的进展，一切时空中的本性都从它获得可感的存在与形式。为看见这善，让我们藉着对我们主耶稣基督的信德纯净自己的心，祂说：\BibleQuote{心里洁净的人是有福的，因为他们要看见天主。}\BibleRef{玛 5:8} 因为为看见这善所需的眼睛，不是我们用来看见那遍布空间的光的那双眼睛，这光一部分在这里，一部分在那里，而非全部在每一处。我们要纯净的视觉与辨别力，是指我们在今生所容许的范围内，看见何为正义、何为虔诚、何为智慧之美的那个能力。看见这些事物的人，将它们看得远超所有空间区域的全部，并发现看见这些事物不需要其知觉在空间距离上的扩展，而是需要藉由非物质的影响而得到增强。\par

% structure: p0099 source=h2 target=section reason=relative-heading-rank; base=h2

\section{第四十三章 结论}

49. 既然这种视像被那些由肉身感觉产生、并被想象所保持和修改的幻想严重阻碍，让我们憎恶这个异端，它由于相信自己的幻想，将天主性体描绘为在空间中，甚至无限空间中延伸与扩散，并切断一侧以给恶腾出空间——无法觉察到恶不是本性，而是反本性；并且用这般可见的外观、形式及部分的一致性来美化这恶本身（这一致性在其各种本性中占主导），以至于无法设想任何缺乏那些美好事物的本性，以至于其中被指摘的种种恶都被埋没在无数美好事物之下。\par

在此让我们结束本论着的这一部分。\Person{摩尼}的其他荒谬之处将在后续内容中，藉天主的许可与帮助，予以揭露。\par

\SourceRecord{Translated by Richard Stothert.；Nicene and Post-Nicene Fathers, First Series；Vol. 4.；Edited by Philip Schaff.；Buffalo, NY: Christian Literature Publishing Co.,；1887.；<http://www.newadvent.org/fathers/1405.htm>.}
